\documentclass[10pt]{beamer} %
\usetheme{CambridgeUS}
% \usecolortheme{seahorse}
\usecolortheme{orchid}
% \usepackage[latin1]{inputenc}
\usefonttheme{professionalfonts}
\usepackage{times}
%\usepackage{tikz}
\usepackage{amsmath}
\usepackage{verbatim}
\usepackage{graphicx}
\usepackage{multicol}
%\usetikzlibrary{arrows,shapes}
\setlength{\tabcolsep}{20pt}
\renewcommand{\arraystretch}{1.5}

\usepackage{amsmath,amssymb,amsfonts}
\usepackage{hyperref}
\usepackage{textcomp}
\usepackage[english]{babel}
\usepackage{times}
\usepackage[T1]{fontenc}

\usepackage{xcolor}
\usepackage{color, colortbl}
\usepackage{multirow}

%\usepackage{beamerthemesplit}
%\usetheme{Warsaw}
\setbeamertemplate{footline}[frame number]
%\usefonttheme{structurebold}
%\usefonttheme[onlymath]{serif}
%\usefonttheme{serif}
\usepackage{pgf,pgfarrows,pgfnodes,pgfautomata,pgfheaps,pgfshade}
\usepackage{amsmath,amssymb}
% \usepackage{inputenc}
%\usepackage[latin1]{inputenc}
\usepackage[english]{babel}
\usepackage{times}
\usepackage{hyperref}
\usepackage{colortbl}
\usepackage{helvet}


\usepackage{tikz}
\usetikzlibrary{arrows,shapes,trees, positioning}
\usepackage{graphics}
\usepackage{color}

\pgfdeclarelayer{background}
\pgfdeclarelayer{foreground}
\pgfsetlayers{background,main,foreground}


%\usepackage[dvipsnames]{xcolor}
%\graphicspath{{images/}}
%\renewcommand{\baselinestretch}{1.5}
% \usepackage[utf8]{inputenc} % allow utf-8 input
% \usepackage[T1]{fontenc}    % use 8-bit T1 fonts
%\usepackage{listings}
%\usepackage{hyperref}       % hyperlinks
%\usepackage{url}            % simple URL typesetting
%\usepackage{booktabs}       % professional-quality tables
%\usepackage{amsfonts}       % blackboard math symbols
%\usepackage{nicefrac}       % compact symbols for 1/2, etc.
%\usepackage{microtype}      % microtypography
%\usepackage{lipsum}
% \usepackage{fancyhdr}       % header
%\usepackage{graphicx}       % graphics
%\usepackage{caption}
%\usepackage{subcaption}
%\usepackage[dvipsnames]{xcolor}
%\usepackage{cite}
%\usepackage{epsfig}
%\usepackage{textcomp}
%\usepackage{bm}
%
%\usepackage{algorithm}
%\usepackage{algorithmic}
%\usepackage{dsfont}
%\usepackage{appendixnumberbeamer}
%Header
% \pagestyle{fancy}
% \thispagestyle{empty}
% \rhead{ \textit{ }} 
% \setlength{\tabcolsep}{15pt}


\newcommand{\defn}{\textcolor{red}{Def:~}}
%\AtBeginEnvironment{frame}{\setcounter{myc}{0}}
\newcounter{myc}[section]%framenumber]
\newcommand{\expl}{\textcolor{blue}{\stepcounter{myc}Example~\arabic{myc}:~}}
\newcommand{\theo}{\textcolor{brown}{Theorem:~}}
\newcommand{\coro}{\textcolor{brown}{Corollary:~}}
\newcommand{\remark}{\textcolor{brown}{Remark:~}}
\newcommand{\sam}{\mathcal{S}}
\newcommand{\sig}{\mathcal{F}}
\newcommand{\R}{\mathbb{R}}
\newcommand{\bg}{\boldsymbol{g}}
\newcommand{\bX}{\boldsymbol{X}}
\newcommand{\bx}{\boldsymbol{x}}
\newcommand{\bY}{\boldsymbol{Y}}
\newcommand{\by}{\boldsymbol{y}}
\newcommand{\bt}{\boldsymbol{t}}
\newcommand{\bu}{\boldsymbol{u}}
\newcommand{\bmu}{\boldsymbol{\mu}}
\newcommand{\mc}{\{X_n:n\geq0\}}


\title[]{Statistical Inference and Multivariate Analysis (MA324)}
\subtitle{\textsc{Lecture Slides}\\ Topic 00 \\
   Introduction and Information
}
\author[MA324]{
}
% \vspace{-5mm} 
% \institute[]{Indian Institute of Technology Guwahati}
\date{}


\begin{document}

% ===================================================================== Frame : Title Page
\thispagestyle{empty}
\maketitle

% ===================================================================== Frame : Syllabus
\begin{frame}
   \frametitle{Syllabus (after mid-semester examination)}
   \begin{itemize}
      \item Confidence Interval: Delta Method, Bootstrap confidence interval.
      \item Simple linear regression: least squares estimation, estimation of
         variance, tests of significance, interval estimation.
      \item Multiple linear regression: least squares estimation, estimation of
         variance, tests of significance, interval estimation,
         multicollinearity, residual analysis, PRESS statistic, detection and
         treatment of outliers, lack of fit.
      \item Multivariate analysis: principle component analysis, factor
         analysis, canonical correlations, cluster analysis.
   \end{itemize}
\end{frame}

% ===================================================================== Frame : Books
\begin{frame}
   \frametitle{Books}
   \begin{itemize}
      \setlength\itemsep{1em}
      \item Text Books
         \begin{itemize}
            \item \textit{Introduction to Linear Regression Analysis} by
               \textit{D.~C.~Montgomery, E.~A.~Peck and G.~G.~Vining}.
            \item \textit{Applied Multivariate Statistical Analysis} by
               \textit{R.~A.~Johnson and D.~W.~Wichern}.
         \end{itemize}
      \item Reference Books
         \begin{itemize}
            \item \textit{Applied Regression Analysis} by \textit{N. R. Draper
               and H. Smith}.
            \item \textit{ Applied Linear Regression} by \textit{S. Weisberg}.
            \item \textit{An Introduction to Multivariate Statistical Analysis}
               by \textit{T. W. Anderson}.
            \item \textit{Applied Multivariate Statistical Analysis} by
               \textit{W. K. Hardle and L. Simar}.
         \end{itemize}
   \end{itemize}
\end{frame}

% ===================================================================== Frame : Grading
\begin{frame}
   \frametitle{Grading Policy}
   \begin{itemize}
      \item Weights in different examinations are as follows:
         \begin{itemize}
            \item Quiz III: 15 (Date: April 09, 2026)
            \item End-semester examination: 30 (Date: May 08, 2026)
         \end{itemize}
      \item For each examination, linear scaling will be used (if needed)
      \item Letter grade based on you performance in 5 exams (3 quizzes, MS, ES)
      \item Relative grading scheme
   \end{itemize}
\end{frame}

% ===================================================================== Frame : Course Website
\begin{frame}
   \frametitle{Course Website}
   \hspace{10mm}\url{https://ayonganguly.github.io/ma324.html}
\end{frame}

% == == == == == == == == == == == == == == == == == == == == 
% Frame: Resource Person
\begin{frame}
   \frametitle{Resource Person}
   \begin{itemize}
      \setlength\itemsep{1em}
      \item Instructor:
         \begin{itemize}
            \item Name: Ayon Ganguly
            \item Office: E1-210, Department of Mathematics
            \item Email: aganguly@iitg.ac.in
         \end{itemize}
      \item Teaching Assistants:
         \begin{itemize}
            \item Hitesh Kumar (hiteshk1004@iitg.ac.in)
            \item Prateek Goel (prateek.goel@iitg.ac.in)
         \end{itemize}
   \end{itemize}
\end{frame}


\end{document}
